\metatitle{Real-rooted polynomials: permutations, set partitions and Stirling permutations}

\metadescription{Examples of Sturm sequences and interlacing from permutation
statistics, set partitions, and Stirling permutations: Eulerian polynomials,
peaks, descents, derangements, and more.}
\metakeywords{real-rooted polynomial, Eulerian polynomial, type B Eulerian
polynomial, derangement polynomial, Stirling permutation, inversion sequence,
Sturm sequence, interlacing polynomial}


\section[realRootedWords]{Permutation statistics, words, and set partitions}

This page collects real-rootedness results for statistics on permutations,
words, and set partitions.  Many examples are
\hyperref[sturmSequenceDefinition]{Sturm sequences}, meaning that
\(P_j\interl P_{j+1}\) for all \(j\).  Others use
\hyperref[interlacingSequenceDefinition]{interlacing sequences},
\hyperref[stablePolynomial]{stable polynomials}, or
\hyperref[polyaFrequency]{total-positivity} methods.  For the broader
context of real-rootedness and unimodality, see
\hyperref[polynomialRootIntro]{unimodality and real-rootedness}.
For descent polynomials of \hyperref[linearExtension]{linear extensions} of
posets, including
\hyperref[generalizedSnakeRealRooted]{generalized snake posets}, see
\hyperref[posetRealRooted]{poset operations preserving real-rootedness}.

The examples are grouped into Eulerian-type descent polynomials, other
permutation statistics, set partitions, and Stirling permutations.  The
subsections are meant as a guide to the methods as much as to the objects:
repeated themes are derivative recurrences, compatible or interlacing
refinements, and multivariate stability.


\subsection[sturmEulerian]{Eulerian-type descent polynomials}


\begin{example*}[Eulerian polynomials, descents in permutations]\label{ex:eulerianSturm}

The \defin{Eulerian polynomials}\label{eulerianPolynomial} are
\[
  A_n(t) \coloneqq \sum_{\sigma \in \symS_n} t^{\des(\sigma)}.
\]
These satisfy
\[
A_n(t) = (1+t(n-1))A_{n-1}(t) + t(1-t) A'_{n-1}(t), \quad A_1(t)=1
\]
and this recurrence implies that \(A_j(t)\interl A_{j+1}(t)\) for all \(j\).
The coefficient \(n-1\) is large enough that this is not a
direct application of \hyperref[interlacingPolynomials]{Wagner's lemma}; the
root positions must be used more carefully.

This result goes back to \name{Frobenius} \cite{Frobenius1910}.

The proof proceeds as follows.  Introduce
$\tilde{A}_n(x) \coloneqq x A_n(x)$.
The recursion then becomes
\[
\tilde{A}_n(x) = x(1-x)\tilde{A}'_{n-1}(x) + nx \tilde{A}_{n-1}(x).
\]
Since \(\tilde{A}'_{n-1}\interl \tilde{A}_{n-1}\), the two summands have
controlled roots.  The term \(x(1-x)\tilde{A}'_{n-1}(x)\) has negative
leading coefficient, while the sum
\[
x(1-x)\tilde{A}'_{n-1}(x) + nx \tilde{A}_{n-1}(x)
\]
has a positive leading coefficient due to $n \geq 2$ in the recursion.
Drawing the roots on the line gives the desired interlacing.
\end{example*}


\begin{example*}[Multi-Eulerian polynomials]
Let $\symS_M$ be the permutations of some multiset
$1^{\mu_1},2^{\mu_2},\dotsc,n^{\mu_n}$ with $N=|\mu|$.
Define
\[
 A_M(x;q) \coloneqq  \sum_{\pi \in \symS_M } x^{\des(\pi)}q^{\maj(\pi)}.
\]
These are the \defin{multi-Eulerian polynomials}\label{multiEulerianPolynomial}.
Then (MacMahon, \cite{MacMahon1960})
\[
 \frac{A_M(x;q)}{\prod_{j=0}^N (1-x q^j)}
 =
 \sum_{m \geq 0}
 \left(\prod_{j=1}^n \qbinom{\mu_j+m}{m}_q \right) x^m.
\]
The polynomials \(A_M(x;1)\) are then real-rooted, see
\cite[Section~2]{Simion1984}.  Ma and Pan \cite[Thm.~1.11]{MaPan2023}
give a proof of this using interlacing polynomials.
For stability in multivariate Eulerian versions, see
\hyperref[ex:stableMultivariateEulerianOperator]{the stable-polynomial page}
and \cite[Eq.~(4.7)]{BrandenHaglundVisontaiWagner2011}.

\name{Danai Deligeorgaki}, \name{Bin Han}, and \name{Liam Solus}
study colored multiset Eulerian polynomials, which simultaneously generalize
MacMahon's multiset Eulerian polynomials and colored Eulerian polynomials
\cite{DeligeorgakiHanSolus2024x}.  They characterize the symmetric cases and
give sufficient conditions for self-interlacing; this implies real-rootedness,
log-concavity, alternating increase, and bi-gamma-positivity.  Their proofs
use multivariate generalizations of MacMahon's identity above.
\end{example*}

\begin{example*}[Log-concavity from triangular-array recurrences]
\name{Umesh Shankar} gives sufficient conditions for log-concavity of rows of
triangular arrays satisfying a super-recurrence generalizing the standard
Eulerian, Stirling, and Lah recurrences \cite{Shankar2025x}.  The proof uses a
weighted-lattice-path interpretation and a weight-increasing injection, and it
also introduces a two-parameter Eulerian analogue with palindromicity and
questions about \hyperref[gammaPositivePolynomial]{gamma-nonnegativity} and
real-rootedness.
\end{example*}

\begin{example*}[Multi-dimensional permutations]
\name{Shaoshi Chen}, \name{Hanqian Fang}, \name{Sergey Kitaev}, and
\name[Candice X. T. Zhang]{Candice X.~T.~Zhang} extend permutation-pattern
statistics to multi-dimensional permutations \cite{ChenFangKitaevZhang2025}.
For canonical three-dimensional permutations, let $R_{n,k}$ count the
permutations with exactly $k$ repeated levels under the max-level statistic, and
set $R_{n,*}(y) \coloneqq \sum_k R_{n,k}y^k$.  They prove that
$\{R_{n,*}(y)\}_{n\geq 0}$ is a generalized Sturm sequence, so each
$R_{n,*}(y)$ is real-rooted.  They also show that weakly increasing
three-dimensional permutations are counted by the Springer numbers.
\end{example*}



\begin{example*}[Eulerian polynomials with descents from prescribed set]
\name{Christos A. Athanasiadis} and
\name{Katerina Kalampogia-Evangelinou} show that for any $T \subseteq [n-1]$,
the polynomials
\[
 A_n^T(x) \coloneqq  \sum_{\substack{\pi \in \symS_n \\ \DES(\pi) \subseteq T}} x^{\des(\pi)}
\]
are real-rooted \cite{AthanasiadisEvangelinou2023xII}.  These are
\defin{restricted Eulerian polynomials}\label{restrictedEulerianPolynomial}.
The proof refines the above polynomials by the value of
$\pi(n)$, and then uses a recursion which preserves an interlacing family of
polynomials.
\end{example*}


\begin{example*}[Inversion-weighted Eulerian polynomials with fixed last entry]

Define the refinement
\[
  A_{n,j}(t,q)
  \coloneqq
  \sum_{\substack{\sigma \in \symS_n \\ \sigma(n)=j}}
  t^{\des(\sigma)}q^{\inv(\sigma)} .
\]
For any $q \gt 0$,
\[
 A_{n,n}(t,q), A_{n,n-1}(t,q), \dotsc, A_{n,1}(t,q)
\]
is an interlacing sequence, see \cite{SavageVisontai2015}.
A short proof is given in \cite{AthanasiadisEvangelinou2023x}.
\end{example*}


\begin{example*}[Type $B$ Eulerian polynomials]\label{ex:typeBEulerianSturm}

Now consider \hyperref[permutationsTypeB]{permutations of type $B$}.
For \(\pi\in B_n\), set \(\pi(0)\coloneqq 0\), and let \(\des(\pi)\) be the cardinality of
\[
 \{ i \in \{0,1,2,\dotsc,n-1\} : \pi(i) \gt \pi(i+1) \}.
\]
The \defin{type $B$ Eulerian polynomials}\label{typeBEulerianPolynomial}
\(P_n(t)=\sum_{\sigma \in B_n} t^{\des(\sigma)}\) then satisfy the recurrence
\begin{equation}
 P_n(t) = ((2n-1)t+1) P_{n-1}(t) + 2t(1-t) P'_{n-1}(t),
\end{equation}
from which the interlacing property follows, see \cite{Chow2022}.
They in fact show that the corresponding operator preserves
real-rootedness, using the method of \name{Julius Borcea} and
\name{Petter Brändén} on \hyperref[stablePolynomial]{stable polynomials}.

The analogous polynomials for type \(D\) are proved to be real-rooted, but
are only conjectured to form interlacing pairs; see \cite[Section~5]{Chow2022}.
\end{example*}




\begin{example*}[Binomial Eulerian polynomials]

The \defin{binomial Eulerian polynomials}\label{binomialEulerianPolynomial}
can be defined as
\[
 \tilde{A}_n(t) \coloneqq \sum_{\sigma \in \symS_n} (1+t)^{\fix(\sigma)} t^{\exc(\sigma)}.
\]
Athanasiadis conjectured that these are real-rooted,
and it was later shown by \name{Jim Haglund} and
\name[P.B. Zhang]{Philip B. Zhang} \cite{HaglundZhang2019} that this holds,
by means of interlacing polynomials.  Their result is stronger: it proves
real-rootedness for a binomial version of the
\hyperref[sEulerianPolynomial]{$\svec$-Eulerian polynomials} defined below,
and the case \(\svec=(2,3,\dotsc,n)\) recovers \(\tilde{A}_n(t)\) via
Lehmer codes.

Computations suggest that $\tilde{A}_n(t) \interl \tilde{A}_{n+1}(t)$, but
this does not follow immediately from the reference above.
\end{example*}


\begin{example*}[Ascents and descents in inversion sequences]
Fix $\svec = (s_1,\dotsc,s_n)$ and define
\[
\mathcal{I}_n^\svec \coloneqq
\{ (e_1,\dotsc,e_n) \in \setZ^n : 0 \leq e_i \lt s_i,\; 1 \leq i \leq n \}
\]
with the conventions that $e_0 = e_{n+1} =0$ and $s_0 = s_{n+1}=1$.
This is the set of
\defin{inversion sequences}\label{inversionSequence} associated with $\svec$.

An index $0 \leq i \leq n$ is an
\defin{ascent}\label{inversionSequenceAscent},
\defin{collision}\label{inversionSequenceCollision}, or
\defin{descent}\label{inversionSequenceDescent}
of an inversion sequence $\evec$ if
\[
 \frac{e_i}{s_i} \lt \frac{e_{i+1}}{s_{i+1}}, \qquad
 \frac{e_i}{s_i} = \frac{e_{i+1}}{s_{i+1}}, \qquad
 \frac{e_i}{s_i} \gt \frac{e_{i+1}}{s_{i+1}},
\]
respectively. Moreover, let $\mathcal{D}_n^\svec \subset \mathcal{I}_n^\svec$
be the set of inversion sequences without collisions.

It was proved in \cite{SavageVisontai2015} and
\cite{GustafssonSolus2020}, respectively, that
\[
E^\svec_n(t) = \sum_{\evec \in \mathcal{I}_n^\svec} t^{\asc(\evec)} \quad \text{and} \quad
D^\svec_n(t) = \sum_{\evec \in \mathcal{D}_n^\svec} t^{\asc(\evec)}
\]
are real-rooted. The methods use the theory of interlacing sequences of polynomials.
In particular, if $\svec$ is a fixed \emph{sequence} of positive integers,
then $E^\svec_n(t) \interl E^\svec_{n+1}(t)$.

The proof by Savage and Visontai also shows that the
\defin{\(\svec\)-Eulerian polynomials}\label{sEulerianPolynomial}
are real-rooted.  See also the proof in
\cite[Sec.~8.2]{Branden2015}.

\name{Jim Haglund} and \name[P.B. Zhang]{Philip B. Zhang}
\cite{HaglundZhang2019} prove a binomial analogue of this theorem.
Define the number of collisions by
\[
  \operatorname{coll}(\evec)
  \coloneqq
  \left|\left\{0\leq i\leq n :
  \frac{e_i}{s_i}=\frac{e_{i+1}}{s_{i+1}}\right\}\right|.
\]
Their \defin{binomial \(\svec\)-Eulerian polynomial}
\label{binomialSEulerianPolynomial} is
\[
  \widetilde{E}_n^\svec(t)
  \coloneqq
  \sum_{\evec\in\mathcal{I}_n^\svec}
  (1+t)^{\operatorname{coll}(\evec)} t^{\asc(\evec)}.
\]
They prove that \(\widetilde{E}_n^\svec(t)\) is real-rooted for every
positive integer sequence \(\svec\).  More precisely, after refining by the
last entry \(e_m=k\), the polynomials
\[
  p_{m,0}^\svec(t),p_{m,1}^\svec(t),\dotsc,p_{m,s_m-1}^\svec(t)
\]
form an interlacing sequence for every \(m\).
The proof writes the recurrence from \(m-1\) to \(m\) as multiplication by a
matrix whose entries are \(1\), \(t\), and \(1+t\).  They then use
\hyperref[realMatrixRecursion]{Brändén's matrix criterion} for interlacing
preservers: it is enough to check the relevant \(2\times2\) submatrices.
The \(1,t\)-only cases are known, and the new \(1+t\) cases are checked by
factoring the small matrices into already known interlacing-preserving
matrices, together with the convexity closure for interlacing sums.
\end{example*}



\subsection[sturmPermutations]{Stirling numbers and pattern statistics}


\begin{example*}[Type 1 Stirling numbers (cycles, left to right minima)]

The unsigned Stirling numbers of the first kind appear in the identity
\[
t(t+1)(t+2)\dotsm (t+n-1) = \sum_{k=1}^n |s(n,k)|t^k
= \sum_{\pi \in \symS_n} t^{\mathrm{numberOfCycles}(\pi)} = \sum_{\pi \in \symS_n} t^{\mathrm{leftToRightMinima}(\pi)}
\]
where $|s(n,k)|$ denotes the unsigned Stirling number of the first kind;
see \oeis{A132393}.
\end{example*}

\begin{example*}[Higher-order Stirling triangles]
\name{Bishal Deb} and \name{Alan D. Sokal} introduce higher-order Stirling
cycle and subset triangles and formulate total-positivity and
\hyperref[realRootedPolynomial]{real-rootedness} questions for them
\cite{DebSokal2025x}.  Their framework asks for total
positivity of the triangle, its row reversal, Toeplitz total positivity of row
sequences, and coefficientwise Hankel total positivity of the row-generating
polynomials.  Toeplitz total positivity is the
\hyperref[polyaFrequency]{Pólya frequency} condition on each row sequence.
They prove several cases for order \(2\).
\end{example*}

\begin{example*}[Big descents in pattern-avoiding permutations]
A descent $i$ of a permutation $\pi$ is called a
\defin{big descent}\label{bigDescent} if
$\pi_i>\pi_{i+1}+1$.  \name{Sergi Elizalde}, \name{Johnny Rivera}, and
\name{Yan Zhuang} classify the big-descent distributions on permutations
avoiding one or two patterns of length three \cite{ElizaldeRiveraZhuang2026}.
Among their future directions, they prove real-rootedness for the
big-descent polynomial on $231$-avoiding permutations and formulate
real-rootedness, log-concavity, and
\hyperref[schurPositivity]{Schur-positivity} questions for the
associated quasisymmetric functions.
\end{example*}


\subsection[peakRunDerangementPolynomials]{Peaks, runs, and derangements}


\begin{example*}[Peak polynomials, peaks in permutations]\label{ex:peakSturm}

The \defin{peak-generating polynomials}\label{peakGeneratingPolynomial}
\[
  P_n(x)=\sum_{\sigma\in\symS_n} x^{\text{peaks}(\sigma)}
\]
have coefficients \oeis{A008303}.
Real-rootedness was proved in \cite{WarrenSeneta1996}, by using the recursion
\begin{equation}
 P_n(x) = (2+x(n-2)) P_{n-1}(x) + 2x(1-x) P'_{n-1}(x).
\end{equation}

After setting \(Q_n=xP_n\), the recursion becomes
\begin{equation}
 Q_n(x) = n x\cdot  Q_{n-1}(x) + 2(1-x)x Q'_{n-1}(x).
\end{equation}

The operator \(T_n:f \mapsto nx f + 2(1-x)x f'\) does not by itself
give a real-rootedness preserver through the stable-operator criterion.
However, the Ma--Wang theorem \cite{MaWang2008} applies to the recurrence,
and gives that all roots of \(Q_n\) are real and that
\(Q_n\interl Q_{n+1}\).
\end{example*}


\begin{example*}[Alternating runs]
\name{Shi-Mei Ma} and \name{Yi Wang} consider the
\defin{alternating-run polynomial}\label{alternatingRunPolynomial}
with coefficients \oeis{A059427}:
\[
  P_n(x) = x \cdot \sum_{\sigma \in \symS_n} x^{\text{peaks}(\sigma)+\text{valleys}(\sigma)}.
\]
The polynomials satisfy the recursion
\[
P_n = x(nx-2x+2)P_{n-1} + x(1-x^2) P'_{n-1}.
\]
Ma and Wang's general theorem gives the
interlacing relation $P_n \interl P_{n+1}$ here as well
\cite{MaWang2008}.  However, note that $x=-1$ is a root with multiplicity
$\lfloor n/2 \rfloor -1$.  See \cite{MaWang2008} for more
references---real-rootedness was observed earlier.
\end{example*}


\begin{example*}[Descents and left peaks in simsun permutations]
Define the polynomials
\[
  P_n(t)=\sum_{\sigma\in \mathrm{RS}_n} t^{\des(\sigma)},
\]
where the sum is over all \hyperref[namedPermutations]{simsun permutations}.
\name{Chak-On Chow} and \name{Wai Chee Shiu} prove that
\begin{equation}
 P_n(t) = ((n-1)t + 1) P_{n-1}(t) + t(1-2t) P'_{n-1}(t), \qquad P_0(t)=1.
\end{equation}
From this recursion, it follows that the \(P_n(t)\) are real-rooted and
interlacing; see \cite[Thm.~2.1 (iii)]{ChowShiu2011}.
They also note that this is the same as counting left peaks in simsun
permutations.

\name{Shi-Mei Ma} and \name{Yeong-Nan Yeh} later prove that interior peaks
also give rise to real-rooted polynomials \cite{MaYeh2016}.
\end{example*}




\begin{example*}[Descents in run-sorted permutations]
In joint work with \name{Olivia Nabawanda}, we show that the
\defin{run-sorted descent polynomials}\label{runSortedDescentPolynomial}
$R_n(t)$ given by $R_1(t)=R_2(t)=t$ and the recursion
\begin{equation}
R_n(t) = t\cdot R'_{n-1}(t) + t\cdot (n-2) R_{n-2}(t),
\end{equation}
satisfy $R_{n-1} \interl R_{n}$ for all $n \geq 1$
\cite{AlexanderssonNabawanda2021}.
The triangle of coefficients can be found in \oeis{A124324}.

The polynomial $R_n(t)$ is the generating function for the number of runs in
\emph{run-sorted permutations}:
\[
 R_n(t) = \sum_{\pi \in \text{RSP}(n)} t^{\des(\pi)+1}.
\]
These are therefore closely related to Eulerian polynomials.
\end{example*}


\begin{example*}[Excedances and derangements, with proof]
The \defin{derangement polynomial}\label{derangementPolynomial} is
$P_n(t)\coloneqq \sum_{\pi \in D(n)} t^{\exc(\pi)}$,
where $D(n)$ is the set of derangements (fixed-point-free permutations) of size $n$,
see \oeis{A046739}. Then $P_{n-1} \interl P_n$.

The proof is as follows.
From \cite{MantaciRakotondrajao2003}, one obtains \(P_1=0\), \(P_2=t\),
and for \(n\geq 3\),
\[
 P_n = t \left( (n - 1) P_{n - 2} + (n-1)  P_{n - 1} + (1-t) P'_{n - 1} \right).
\]
Note that $P_n(t)$ is a polynomial of degree $n-1$, as the
permutation $\pi=2345\dotsc n 1$ contributes with $t^{n-1}$ in the definition.

Proceed by induction over \(n\), and assume that
\(P_{n-2}\interl P_{n-1}\) for some \(n\geq 3\).

By the induction hypothesis, and using that $P'_{n-1} \interl P_{n-1}$, we can deduce that
\[
  P_{n-2} \interl P_{n-1}, \text{ and }
 (n-1)P_{n-1} + (1-t)P'_{n-1} \interl P_{n-1}.
\]
Using \cite{Wagner1992}, we have
\begin{align*}
 (n-1)P_{n - 2} + (n-1)P_{n - 1} + (1-t)P'_{n-1} &\interl P_{n-1} \\
 t\left((n-1)P_{n - 2} + (n-1)P_{n - 1} + (1-t)P'_{n-1}\right) &\interl t P_{n-1} \\
 P_{n} &\interl t P_{n-1}.
\end{align*}
It follows that \(P_{n-1}\interl P_n\).

See also \cite{LiuWang2007} where this is proved using a different recursion.
\end{example*}


\section[sturmSetPartitions]{Set partitions}



\begin{example*}[Touchard polynomials (blocks in set partitions), with proof]\label{touchardPolynomial}

The \defin{Touchard polynomials}\label{touchardPolynomialDefinition}
may be defined via
\[
T_n(x) \coloneqq \sum_{P \in SP(n)} x^{\mathrm{blocks}(P)}
= \sum_{k=1}^n S(n,k) x^k,
\]
where $SP(n)$ is the set of set partitions of $[n]$, and we count the number
of blocks.
Note that the coefficients are Stirling numbers of the second type.

Wilf gives the recurrence \cite[pp.~20, 139]{Wilf1994}
\[
  T_n(x) = x \cdot  T_{n-1}(x) + x\cdot T'_{n-1}(x).
\]
We know $T_{n-1} \interl x T_{n-1}$. Moreover, $T_{n-1} \interl x T'_{n-1}$,
so by \cite{Wagner1992}, we have $T_{n-1} \interl x T_{n-1} +  x T'_{n-1}$.
But this is exactly $T_{n-1} \interl T_n$.

Another way to prove real-rootedness is by interpreting \(T_n(x)\) as a
\hyperref[touchardAsMatching]{matching polynomial}.
\end{example*}


\begin{example*}[Set partitions without singleton blocks]

Define \(D_n(t)=\sum_{\pi \in SP_2(n)} t^{\mathrm{blocks}(\pi)}\),
where $SP_2(n)$ is the set of set partitions of $[n]$ whose blocks all have
size at least $2$.
Then
\[
 D_n(t) = t\left( D'_{n-1}(t) + (n-1) D_{n-2}(t) \right)
\]
and $D_{n-1}(t)$ and $D_n(t)$ have interlacing roots, see
\cite[Thm.~1]{BonaHezo2016}.
Note that the degree of $D_n(t)$ is $\lfloor n/2 \rfloor$.


The proof is short by using Wagner's lemma and induction.
The base case is immediate; assume $D_{n-2} \interl D_{n-1}$ for some
$n \geq 2$.
We then know that
\[
D'_{n-1} \interl D_{n-1} \text{ and } (n-1)D_{n-2} \interl D_{n-1}.
\]
Wagner's lemma now gives that
\[
D'_{n-1} + (n-1)D_{n-2} \interl D_{n-1}
\implies
D_{n-1} \interl t \left( D'_{n-1} + (n-1)D_{n-2} \right),
\]
and this is precisely $D_{n-1} \interl D_n$, which completes the induction step.

This result also follows from the general machinery in
\cite[Subsec.~3.4]{Hitczenko2024x}.
A recursion for set partitions into blocks of size at least $s$ is also mentioned
there (originally due to \name{Louis Comtet})
but it does not immediately imply real-rootedness.

\end{example*}



\begin{example*}[Type $B$ set partitions]

\name{David Wang} considers a type $B$ version, and more generally a colored
version, of counting blocks in set partitions \cite{Wang2014}.
The polynomials considered satisfy the (general) recursion
\[
  T_n(x) = (x+c) \cdot  T_{n-1}(x) + mx\cdot T'_{n-1}(x),
\]
where $m$ and $c$ are positive integers.

\end{example*}


\begin{example*}[Ordered partitions]
Simion considers any multiset $M$ and the number $O(M,k)$ of compositions into
exactly $k$ parts \cite{Simion1984}.  The polynomial
$\sum_k O(M,k)x^k$ is real-rooted.

In particular, this implies the real-rootedness of the Touchard polynomials.
\end{example*}


\subsection[sturmOther]{Stirling and quasi-Stirling permutations}


\begin{example*}[Stirling permutations]

The descent number of a \hyperref[permutationsStirling]{Stirling permutation}
\(\pi \in Q_n\) is the cardinality of
\[
 \{ i : \pi_i \gt \pi_{i+1} \} \cup \{2n\}.
\]
The \defin{Stirling permutation descent polynomials}
\label{stirlingPermutationDescentPolynomial} satisfy the recurrence
\begin{equation}
 P_n(t) = (2n+1)t \cdot P_{n-1}(t) + t(1-t) P'_{n-1}(t),
\end{equation}
see \cite[Eq.~13]{Carlitz1965} and \cite[Thm.~2.1]{GesselStanley1978}.
The coefficients can be found in \oeis{A008517}; the first few values are
presented below.
\begin{array}{rrrrrr}
\toprule
 1  \\
 1 & 2  \\
 1 & 8 & 6 \\
 1 & 22 & 58 & 24 \\
 1 & 52 & 328 & 444 & 120 \\
 1 & 114 & 1452 & 4400 & 3708 & 720 \\
\bottomrule
\end{array}
These numbers appear in the Hilbert series of the Whitehouse complex;
see \cite{Readdy2005}.  These polynomials are the
\hyperref[hStarPolynomial]{$h^*$-polynomials} of the
\hyperref[orderPolytope]{order polytope} defined by the inequalities
\[
 y_1 \leq y_2 \leq \dotsb \leq y_n \qquad 0 \leq y_j \leq x_j \leq 1 \text{ for } j=1,\dotsc,n.
\]
\name{Shi-Mei Ma}, \name{Jun Ma}, \name{Jean Yeh}, and
\name{Yeong-Nan Yeh} prove $e$-positivity for multivariate $k$th-order
Eulerian polynomials attached to $k$-Stirling permutations
\cite{MaMaYehYeh2021x}.  They also give tree interpretations for the
coefficients in related gamma-type expansions.
For $n=20$, the linear term of the Ehrhart polynomial (of degree $40$) has
negative coefficient $-22354369153/6983776800$.

\name{Miklós Bóna} proved that these have only real zeros \cite{Bona2009}.
The Ma--Wang theorem also gives real-rootedness and interlacing.

\name{Jim Haglund} and \name{Mirko Visontai} \cite[Thm.~3.3]{HaglundVisontai2012}
show that a multivariate generalization is stable.
\end{example*}

\begin{theorem*}[Dey's interlacing criterion, \cite{Dey2023InterlacingZeroes}]
\label{deyInterlacingCriterion}
Let $(F_n(x))_{n\geq 0}$ be a sequence of polynomials satisfying
\[
  F_n(x)=\bigl(\alpha_1(n)x+\alpha_2\bigr)F_{n-1}(x)
  +\alpha_3x(1-x)F'_{n-1}(x),
\]
where $\alpha_2,\alpha_3\in\setN\cup\{0\}$. Suppose that each
$F_n(x)$ is real-rooted with distinct nonpositive zeros, and that
some integer $k$ satisfies, for all $n\geq k$,
\[
  2\alpha_1(n)-n\alpha_3>0,
  \qquad
  n\alpha_1(n)\alpha_3+3n\alpha_2\alpha_3
  -4\alpha_1(n)\alpha_2>0.
\]
Then $F_{n-1}(x)\interl F_n(x)$ for all $n\geq k$.
\end{theorem*}

Dey applies this criterion to the type $A$ and type $B$ Eulerian
polynomials, to the Stirling permutation polynomials above, and to
\hyperref[restrictedEulerianPolynomial]{restricted Eulerian polynomials}.
Thus it is a useful way to upgrade an
already known real-rootedness proof to a consecutive interlacing
statement when the recurrence has the form above.


\begin{example*}[Quasi-Stirling permutations]
A \defin{quasi-Stirling permutation}\label{quasiStirlingPermutation}
of size $n$ is a permutation of the multiset
$\{1,1,2,2,\dotsc,n,n\}$ avoiding the patterns $1212$ and $2121$.
Let $\widetilde{Q}_n$ be the set of quasi-Stirling permutations of size $n$,
and set
\[
 \widetilde{Q}_n(t) \coloneqq \sum_{\pi \in \widetilde{Q}_n} t^{\des(\pi)},
\]
where the final position is counted as a descent, as above.

\name{Sergi Elizalde} \cite[Thm.~3.2]{Elizalde2021} proved that
$\widetilde{Q}_n(t)$ has real, distinct, and nonpositive roots.
More generally, for $1 \leq r \leq n$, let
\[
 A_{n,r}(t) \coloneqq \sum_{\pi \in \symS_n} t^{\exc_r(\pi)}, \qquad
 J_{n,r}(t) \coloneqq \sum_{\pi \in \operatorname{Inj}([n-r],[n])} t^{\exc(\pi)},
\]
where $\exc_r(\pi) \coloneqq |\{i:\pi_i \geq i+r\}|$.
Then $A_{n,r}(t)$ and $J_{n,r}(t)$ also have real, distinct, and nonpositive
roots.
The quasi-Stirling case follows from the relation
\[
 \widetilde{Q}_n(t)=t J_{2n,n+1}(t).
\]
Elizalde also proves asymptotic normality for the descent distribution and
extends the analysis to a one-parameter family generalizing
\(k\)-Stirling permutations.
\end{example*}
